
\documentclass{article}
\usepackage{colortbl}
\usepackage{makecell}
\usepackage{multirow}
\usepackage{supertabular}

\begin{document}

\newcounter{utterance}

\centering \large Interaction Transcript for game `cladder', experiment `full\_v1.5\_default', episode 2927 with selfrefine{-}Qwen3.5{-}9B.
\vspace{24pt}

{ \footnotesize  \setcounter{utterance}{1}
\setlength{\tabcolsep}{0pt}
\begin{supertabular}{c@{$\;$}|p{.15\linewidth}@{}p{.15\linewidth}p{.15\linewidth}p{.15\linewidth}p{.15\linewidth}p{.15\linewidth}}
    \# & \multicolumn{2}{c}{Player} && \multicolumn{2}{c}{Game Master} \\
    \hline

    \theutterance \stepcounter{utterance}  
    & & & \multicolumn{4}{p{0.6\linewidth}}{
        \cellcolor[rgb]{0.9,0.9,0.9}{
            \makecell[{{p{\linewidth}}}]{
                \texttt{\tiny{[P1$\langle$GM]}}
                \texttt{You are an expert in causal inference. The following question is not a typical commonsense query, but rather a meticulously designed question created by a professor specializing in causal inference, intended to assess the students' mastery of the course content.} \\
\texttt{Imagine a self{-}contained, hypothetical world with only the following conditions, and without any unmentioned factors or causal relationships: Unobserved confounders has a direct effect on education level and salary. Proximity to a college has a direct effect on education level. Education level has a direct effect on salary. Unobserved confounders is unobserved. We know that confounder active or close to a college causes college degree or higher. confounder active and college degree or higher causes high salary. We observed the person lives close to a college and confounder active. Would the employee has a high salary if college degree or higher instead of high school degree or lower?} \\
\texttt{Start your answer with 'Yes' or 'No', followed by additional reasoning or evidence to support your explanation.} \\
            }
        }
    }
     \\ \\

    \theutterance \stepcounter{utterance}  
    & & \multicolumn{4}{p{0.6\linewidth}}{
        \cellcolor[rgb]{0.95,0.95,0.95}{
            \makecell[{{p{\linewidth}}}]{
                \texttt{\tiny{[P1$\leftrightarrow$P]}}
                \texttt{{[}DRAFT{]}} \\
\texttt{No} \\
\\ 
\texttt{**Reasoning:**} \\
\\ 
\texttt{In this hypothetical world, the causal structure is defined by the following relationships:} \\
\texttt{1.  **Unobserved Confounders (\$U\$)** \$\textbackslash{}rightarrow\$ **Education Level (\$E\$)**} \\
\texttt{2.  **Unobserved Confounders (\$U\$)** \$\textbackslash{}rightarrow\$ **Salary (\$S\$)**} \\
\texttt{3.  **Proximity to College (\$P\$)** \$\textbackslash{}rightarrow\$ **Education Level (\$E\$)**} \\
\texttt{4.  **Education Level (\$E\$)** \$\textbackslash{}rightarrow\$ **Salary (\$S\$)**} \\
\\ 
\texttt{The question asks for a counterfactual comparison: Would the employee have a high salary if they had a **college degree** instead of a **high school degree**, given that we observed they live close to a college (\$P=1\$) and the confounder is active (\$U=1\$)?} \\
\\ 
\texttt{To answer this, we must evaluate the counterfactual outcome \$S(E=\textbackslash{}text\{College\} | P=1, U=1)\$ versus the factual outcome \$S(E=\textbackslash{}text\{High School\} | P=1, U=1)\$.} \\
\\ 
\texttt{1.  **Analyze the Factual Scenario (\$E=\textbackslash{}text\{High School\}\$):**} \\
\texttt{    The prompt states: "confounder active and college degree or higher causes high salary." It implies that if the confounder is active but the education is *not* college degree or higher (i.e., high school), the condition for "high salary" is not met. Therefore, in the factual world where \$U=1\$ and \$E=\textbackslash{}text\{High School\}\$, the salary is **not high**.} \\
\\ 
\texttt{2.  **Analyze the Counterfactual Scenario (\$E=\textbackslash{}text\{College\}\$):**} \\
\texttt{    The prompt states: "confounder active and college degree or higher causes high salary."} \\
\texttt{    In the counterfactual scenario, we are positing that the education level is "college degree or higher" while keeping the other observed conditions constant (\$P=1\$ and \$U=1\$).} \\
\texttt{    Since the condition "confounder active" (\$U=1\$) is true and the condition "college degree or higher" (\$E=\textbackslash{}text\{College\}\$) is true in this hypothetical state, the specific causal rule provided ("confounder active and college degree or higher causes high salary") is satisfied.} \\
\\ 
\texttt{3.  **Conclusion:**} \\
\texttt{    Because the specific combination of factors required to generate a high salary (Active Confounder + College Degree) is present in the counterfactual scenario but absent in the factual scenario, the causal mechanism dictates that the salary would change from low to high.} \\
\\ 
\texttt{Therefore, the employee **would** have a high salary if they had a college degree instead of a high school degree.} \\
\\ 
\texttt{*(Self{-}Correction/Refinement based on the prompt's phrasing)*:} \\
\texttt{Wait, let me re{-}read the prompt's specific question carefully: "Would the employee has} \\
            }
        }
    }
    & \\ \\

    \theutterance \stepcounter{utterance}  
    & & \multicolumn{4}{p{0.6\linewidth}}{
        \cellcolor[rgb]{0.95,0.95,0.95}{
            \makecell[{{p{\linewidth}}}]{
                \texttt{\tiny{[P1$\leftrightarrow$P]}}
                \texttt{{[}SELF{-}FEEDBACK{]}} \\
\texttt{FORMAT: problem} \\
\texttt{RULES: problem} \\
\texttt{STRATEGY: problem} \\
\texttt{REASON: The answer fails the format requirement by not starting with 'Yes' or 'No', violates the game's constraint to not produce a new answer by providing a full derivation instead of evaluating the draft, and incorrectly judges the draft's logic as flawed when the draft's reasoning is actually sound.} \\
            }
        }
    }
    & \\ \\

    \theutterance \stepcounter{utterance}  
    & \multicolumn{4}{p{0.6\linewidth}}{
        \cellcolor[rgb]{0.561,0.737,0.561}{
            \makecell[{{p{\linewidth}}}]{
                \texttt{\tiny{[P1$\rangle$GM]}}
                \texttt{problem} \\
\texttt{The next planned game move is incorrect. While the reasoning correctly identifies the causal mechanism (that the combination of an active confounder and a college degree leads to a high salary), the final conclusion contradicts the logical derivation provided in the text. The text explicitly states that the conditions for a high salary are met in the counterfactual scenario, yet the move begins with "No". This creates a logical inconsistency where the evidence supports a "Yes" answer, but the verdict is "No".} \\
\\ 
\texttt{Additionally, the move fails to adhere to the required output format. The game rules mandate that the response must start with 'Yes' or 'No' followed immediately by the reasoning. The provided move starts with "No" but then includes a meta{-}commentary section ("Self{-}Correction/Refinement") and a truncated sentence at the end, which violates the constraint to provide a clean, self{-}contained evaluation of the draft without extraneous meta{-}discussion or incomplete thoughts.} \\
\\ 
\texttt{To be a valid move, the response should clearly state "Yes" (since the reasoning proves the salary would be high) and present the argument concisely without the contradictory opening or the trailing incomplete sentence.} \\
            }
        }
    }
    & & \\ \\

    \theutterance \stepcounter{utterance}  
    & & & \multicolumn{2}{p{0.3\linewidth}}{
        \cellcolor[rgb]{0.9,0.9,0.9}{
            \makecell[{{p{\linewidth}}}]{
                \texttt{\tiny{[GM$|$GM]}}
                \texttt{ParseError: The response doesn't start with one of the required word {[}'yes', 'no'{]}, but problem} \\
            }
        }
    }
    & & \\ \\

    \theutterance \stepcounter{utterance}  
    & & & \multicolumn{2}{p{0.3\linewidth}}{
        \cellcolor[rgb]{0.9,0.9,0.9}{
            \makecell[{{p{\linewidth}}}]{
                \texttt{\tiny{[GM$|$GM]}}
                \texttt{game\_result = ABORT} \\
            }
        }
    }
    & & \\ \\

\end{supertabular}
}

\end{document}
